% 4.4(c) patch -- direct-kernel ranking control over decoder width k
% POST-HOC SUPPLEMENTARY ANALYSIS; NOT PREREGISTERED; DEVELOPMENT SEEDS ONLY.
% Frozen holdout seeds 301-305 were NOT re-run.
\paragraph{Post-hoc supplementary analysis.}
This analysis is post-hoc and supplementary, was not preregistered, uses development seeds 201-205 only, and was not used to modify any main-experiment hyper-parameter. Frozen holdout seeds 301-305 were not re-run; the $k=5$ holdout values serve as an independent reference point only.

\paragraph{Scope.}
Holding the frozen amount-free renormalised cost matrix, $\epsilon = 0.05$ and
$\lambda = 0.5$ fixed, we sweep the mutual-selection width
$k \in \{2, 3, 5, 7, 10, 15\}$ and ask how the relative ordering of
\textsc{Conditional} and \textsc{AmountFree} (direct-kernel ranking) moves on the
\emph{development} distribution (seeds 201--205). The sweep answers how the relative
ordering changes with $k$ on development; it does \emph{not} estimate holdout performance
for $k \neq 5$, because seeds 301--305 are never re-run and the frozen holdout exists only
at the default $k=5$.

\paragraph{Result.}
In the post-hoc k sweep restricted to development seeds 201-205, the conditional-minus-direct-kernel difference is -0.0371 at $k=3$ with a two-sided paired permutation $p = 0.0002$, triggering the pre-specified contribution-downgrade condition C. We therefore no longer describe conditional decoding as generally superior to direct-kernel ordering: its role narrows to repairing the ordering of the raw UOT plan, and its relative merit against the direct kernel is configuration- and split-dependent.

\begin{table}[t]
\centering
\small
\begin{tabular}{rrrrrrr}
\hline
$k$ & RAW & COND & AMOUNT\_FREE & SUPPORT$+$K & COND $-$ AMT & paired $p$ \\
\hline
2 & 0.0008 & 0.3504 & 0.3861 & 0.3786 & -0.0357 & 0.0002 \\
3 & 0.1236 & 0.4797 & 0.5169 & 0.5067 & -0.0371 & 0.0002 \\
5 & 0.2374 & 0.3129 & 0.3172 & 0.3112 & -0.0043 & 0.0571 \\
7 & 0.2623 & 0.2820 & 0.2828 & 0.2775 & -0.0008 & 0.4598 \\
10 & 0.2594 & 0.2720 & 0.2621 & 0.2571 & +0.0099 & 0.0006 \\
15 & 0.2517 & 0.2589 & 0.2329 & 0.2303 & +0.0260 & 0.0003 \\
\hline
\end{tabular}
\caption{Macro-edge F1 by decoder width $k$ (development seeds 201--205, unweighted mean
over the 15 bridge--seed cells; $\pm$1 SD over the five seeds per bridge is shown in the
figure). The paired column is a two-sided paired sign-flip permutation test over the 15
paired bridge--seed cells ($n_{\mathrm{perm}} = 20000$, RNG seed 20240101).
With 15 pairs the smallest attainable two-sided exact $p$-value is $6.1035\times10^{-5}$;
with only five pairs per bridge it is $0.0625$, so per-bridge five-cell tests must not be
reported as $p<0.05$ findings. $k=3$ is the single pre-specified confirmatory trigger; all
other $p$-values are descriptive.}
\end{table}

\paragraph{Development vs.\ frozen holdout at $k=5$.}
The development difference at $k=5$ is $-0.0043$ while the pre-existing frozen
holdout value of the same contrast is $+0.0048$: the sign reverses. Neither
direction can therefore be summarised as a general advantage across data splits.

\paragraph{Hyper-parameters.}
The paper default $k=5$ is unchanged. No hyper-parameter was re-selected from this
experiment, and no result here replaces any main-experiment result.
